\section{Evaluation}
\label{sec:evaluation}

The verified contracts establish numerical guarantees on normal return.  We
evaluate three complementary empirical questions: the
tightness of the proved error bounds, the practical significance of calls
rejected by the runtime checks, and the performance cost of the
wrappers.  We investigate the following research questions:

\begin{description}
  \item[(RQ1) Usefulness of the error bounds.]
    For results produced by accepted calls to the default wrappers for the six
    double-word operations, how closely do the observed errors approach the
    proved bounds?
    As secondary context, how do the errors produced by QD compare on the same
    inputs?

  \item[(RQ2) Usefulness of the wrappers.]
    For generated inputs that cause a runtime check in each wrapper to
    fail, how often does QD return a finite, valid double word whose
    MPFR-measured error is no greater than the corresponding \FenceDD{} error
    bound?

  \item[(RQ3) Cost of the wrappers.]
    On a common set of inputs for which every evaluated call is accepted, what
    runtime overhead does each wrapper add relative to the
    corresponding unchecked kernel, and how does its end-to-end runtime
    compare with QD?
\end{description}

\subsection{Experimental Setup}
\label{sec:evaluation-setup}

We compare \FenceDD{} with QD \citep{Hida2001,QD}, using MPFR to compute
reference values at greater precision \citep{Fousse2007}.  We evaluate both
wrappers for each of the six double-word operations and the single public
wrapper for each error-free transformation.  For every test, both libraries
receive the same binary64 operands, including the same high and low components
of each double-word operand.

To exercise the binary64 range, we generate the same number of inputs at each
base-two exponent.  Binary64 operands cover exponents from $-1074$ through
$1023$, while double-word inputs cover exponents from $-969$ through $1022$.
For 95\% of the double-word inputs, MPFR constructs a 106-bit source value,
which is rounded into high and low components.  The remaining 5\% have a
binary64 high component and a zero low component.  The narrower exponent range
keeps all 106 source bits representable while ensuring that the high component
is finite.  Signs are balanced, and every generated double-word operand is
valid according to \Cref{def:double-word}.

We classify a call as accepted when all runtime checks succeed and the
function returns a result.  A call is rejected when a runtime check stops its
execution.  RQ1 and RQ2 use \RQOneSamplesPerExponent{} inputs at each exponent,
giving 79,680 inputs per double-word operation and 83,920 per transformation.
All floating-point computations use rounding to nearest with ties to even and
gradual underflow.  \Cref{app:evaluation-details} gives the complete input
construction, reference computation, boundary tests, and timing protocol.

\subsection{Usefulness of the Error Bounds (RQ1)}

To address RQ1, we evaluate the six double-word operations on generated inputs
for which calls to the default wrapper are accepted.  MPFR reconstructs the real
values represented by the input and output pairs and computes the
corresponding reference operation at 2112-bit precision.  If $r$ is the MPFR
reference and $\widehat r$ is the value represented by the returned pair, the
observed error is $|\widehat r-r|$.

The bounds in \Cref{tab:operation-preview} have the form
\begin{equation}
  \begin{aligned}
    B(r) &= C u^2|r|+K\eta,\\
    B_{\mathrm{rel}}(r) &= C u^2|r|,
    & B_{\mathrm{abs}} &= K\eta.
  \end{aligned}
  \label{eq:rq1-error-bound}
\end{equation}
We refer to $B_{\mathrm{rel}}(r)$ and $B_{\mathrm{abs}}$ as the relative and
absolute terms.  An input lies in the relative region
when $B_{\mathrm{rel}}(r)\geq B_{\mathrm{abs}}$ and in the absolute region
otherwise.  The addition bounds have no absolute term, so all nonzero addition
results lie in the relative region.

For every result with $B(r)>0$, we measure how closely its error approaches
the bound using
\begin{equation}
  \rho=\frac{|\widehat r-r|}{B(r)}.
  \label{eq:rq1-error-ratio}
\end{equation}
A value of $\rho$ near one indicates that the observed error approaches the
proved bound.  Smaller values indicate that the bound is less tightly attained
on that input.  The numerical theorem implies $\rho\leq1$ for every
\FenceDD{} result produced by an accepted call.  Addition results for which $B(r)=0$ are
reported separately and must be exact.

The QD routines may use a different sequence of binary64 operations from the
corresponding \FenceDD{} kernels.  To provide a common numerical scale, we
divide each finite QD error by the corresponding \FenceDD{} error bound
$B(r)$.  This bound is not proved for QD, and the QD and \FenceDD{} ratios are
not expected to be equal.  Nonfinite QD results are reported separately.

\begin{table}[t]
  \caption{RQ1: Mean and maximum values of the error ratio $\rho$ from
  \Cref{eq:rq1-error-ratio}, separated into the regions determined by the terms
  of the \FenceDD{} error bound in \Cref{eq:rq1-error-bound}.  Region share is
  the percentage of results produced by accepted \FenceDD{} calls that fall in
  each region; the shares for an operation sum to 100\%.  For QD,
  the value of the \FenceDD{} error bound provides a common scale; the bound is
  not proved for QD.  Region shares are reported to two decimal places and
  ratios to four decimal places.}
  \label{tab:rq1-tightness}
  \centering
  \setlength{\tabcolsep}{2pt}
  \begin{tabular}{@{}llllll@{}}
    \toprule
    & & \multicolumn{2}{c}{\FenceDD{} $\rho$}
        & \multicolumn{2}{c}{QD $\rho$} \\
    \cmidrule(lr){3-4}\cmidrule(l){5-6}
    Operation & Share
              & Mean
              & Max
              & Mean
              & Max \\
    \midrule
    \multicolumn{6}{@{}l}{\emph{Relative region}} \\[2pt]
    \input{generated/rq1_relative_compact_rows}\\
    \addlinespace
    \multicolumn{6}{@{}l}{\emph{Absolute region}} \\[2pt]
    \input{generated/rq1_absolute_compact_rows}\\
    \bottomrule
  \end{tabular}
\end{table}

\paragraph{Answer to RQ1.}
\Cref{tab:rq1-tightness} shows that the proved bounds can be approached, but
are conservative for most generated inputs.  Across
\RQOneNormalReturnCount{} results produced by accepted calls, no \FenceDD{}
error exceeded its proved bound.  The largest observed values of $\rho$ were
\RQOneRelativeMaximumRatio{} for \RQOneRelativeMaximumOperation{} in the
relative region and \RQOneAbsoluteMaximumRatio{} for
\RQOneAbsoluteMaximumOperation{} in the absolute region.  By contrast, the
mean for every operation was at most \RQOneRelativeMaximumMean{} in the
relative region.  For every operation with an absolute term, the mean in the
absolute region was at most \RQOneAbsoluteMaximumMean{}.  The absolute region
contained 14--16\% of the results for these operations, so these measurements
represent a substantial portion of the generated tests rather than isolated
boundary cases.
Thus, some observed errors came close to the proved limits, while the bounds
retained substantial margin on average.

As secondary context, QD produced the same observed ratios for both addition
operations and \DWDivFP{}.  Its relative errors were somewhat larger for the
two multiplication operations and smaller for \textsc{DWDivDW}.  One
\DWTimesFP{} result reached $\rho=1.038$, exceeding $B(r)$ by
\RQOneQDExcessPercent{}.  QD may
compute the same mathematical operation using a different algorithm, so
$B(r)$ is not a proved bound for its result.  We use $B(r)$ only to place the
observed QD and \FenceDD{} errors on a common scale.

\FloatBarrier
\subsection{Usefulness of the Wrappers (RQ2)}

A call is rejected when a runtime check cannot establish a condition required
by its verified contract.  To assess the practical significance of these
rejections, we run the corresponding QD operation on every generated input
associated with a rejected \FenceDD{} call and compare its result with an MPFR
reference.

For the six double-word operations, we consider separately calls rejected by
runtime checks in the default wrappers and calls rejected by runtime checks in
the strong wrappers.  These populations overlap: every call rejected in a
default wrapper is also rejected in the corresponding strong wrapper.
Percentages are therefore calculated separately using the number of rejected
calls in the indicated population.  Because the generated tests contain the
same number of inputs at each exponent, these percentages describe our tests
and do not estimate rejection rates in application workloads.

For a double-word operation, the numerical criterion requires QD to return a
finite, valid double word whose error is no greater than the corresponding
\FenceDD{} bound $B(r)$ from \Cref{eq:rq1-error-bound}.  The bound is not proved
for QD; we use it only as a common numerical criterion.  No finite QD result in
these populations was invalid or exceeded $B(r)$, so the table omits those
zero categories.

\begin{table}[t]
  \caption{RQ2: QD results for inputs on which a runtime check stopped a call to
  the indicated wrapper.  Rejected calls gives the population size; the other
  entries are percentages of that population.  The default and strong
  populations overlap.  No finite QD result was invalid or exceeded $B(r)$, so
  those zero categories are omitted.}
  \label{tab:rq2-double-word}
  \centering
  \setlength{\tabcolsep}{3pt}
  \begin{tabular}{@{}llrrr@{}}
    \toprule
    Operation & Wrapper & \shortstack{Rejected\\calls}
              & \shortstack{QD\\nonfinite}
              & \shortstack{Valid,\\error $\leq B(r)$} \\
    \midrule
    \input{generated/rq2_double_word_compact_rows}\\
    \bottomrule
\end{tabular}
\end{table}

Calls stopped by runtime checks usually occurred on inputs for which QD also
failed the numerical criterion.  QD returned a finite, valid double word with
error no greater than the corresponding \FenceDD{} bound for only
\RQTwoDefaultQDWithinPercent{} of calls stopped in the default wrappers and
\RQTwoStrongQDWithinPercent{} of calls stopped in the strong wrappers.

The operation-level results expose important differences.  Every QD result
associated with a rejected call to a default multiplication wrapper was
nonfinite.  Among calls rejected by runtime checks in the default wrappers, QD
met the numerical criterion only for division.  By contrast, QD met the
criterion for every call rejected by the strong addition checks, showing that
these sufficient checks can be conservative.

Each error-free transformation has one public wrapper; it is not classified as
default or strong.  For inputs associated with rejected calls, we compare
\TwoSum{}, \FastTwoSum{}, and \FastTwoMult{} with QD's
\texttt{two\_sum}, \texttt{quick\_two\_sum}, and \texttt{two\_prod},
respectively.  Because the verified contracts require exact reconstruction, a
finite, valid QD pair either reconstructs the MPFR result exactly or fails the
criterion.

\begin{table}[t]
  \caption{RQ2: QD results for inputs on which a runtime check stopped a call to
  the public wrapper for each error-free transformation.  Rejected calls gives
  the population size; the other entries are percentages of that population.}
  \label{tab:rq2-transformations}
  \centering
  \setlength{\tabcolsep}{2pt}
  \begin{tabular}{@{}lrrrrr@{}}
    \toprule
    Transformation & \shortstack{Rejected\\calls}
                   & Nonfinite
                   & Invalid
                   & \shortstack{Not\\exact}
                   & Exact \\
    \midrule
    \input{generated/rq2_eft_compact_rows}\\
    \bottomrule
  \end{tabular}
\end{table}

The results reinforce the pattern in \Cref{tab:rq2-double-word}.  No generated
call to the \TwoSum{} wrapper was rejected.  On inputs associated with rejected
calls to the wrappers for
\FastTwoSum{} and \FastTwoMult{}, QD reconstructed the exact result in only
\RQTwoFastSumExactPercent{} and \RQTwoFastMultExactPercent{} of cases,
respectively.

\paragraph{Where the two double-word wrappers differ.}
For \RQTwoDefaultOnly{} inputs, every runtime check in the default wrapper
succeeded, while the strong check in the corresponding strong wrapper failed.
For multiplication, these cases had a leading-product exponent of $1023$; for
division, they had leading-quotient exponents of $1022$ or $1023$.  The
addition wrappers differed on 196 inputs because the strong checks bound the
operands separately even though the computed sum remained finite.  Checking
result finiteness after kernel execution therefore allows calls to the default
wrapper to be accepted when the corresponding calls to the strong wrapper are
rejected.

\paragraph{Answer to RQ2.}
On these generated tests, runtime checks usually stopped calls on inputs for
which QD also failed the corresponding numerical criterion.  The
operation-level results nevertheless confirm that the checks are sufficient
rather than necessary, particularly for addition.  Because QD may use
different algorithms, these measurements do not establish how an unchecked
\FenceDD{} kernel would behave on the same inputs.

\FloatBarrier
\subsection{Performance Cost of the Wrappers (RQ3)}

To address RQ3, we time the default and strong wrappers for the six
double-word operations and the wrapper for each of the three transformations.
Each timing test uses \RQThreeInputCount{} inputs
whose operand exponents lie between $-100$ and $100$, together with
\RQThreeTrialCount{} trials.  All runtime checks succeed for these inputs.
Before timing, we confirm that every wrapper produces the same
result as its unchecked kernel.

For the comparison with the kernels, each wrapper and its kernel are compiled
with CompCert.  For the comparison with QD, both libraries are compiled with
Clang.  The latter build is used only for the performance comparison and is
not covered by the CompCert correctness theorem.  The experiments run on an
Apple M4 Max with 16 cores.  Compiler versions, options, and the complete
timing procedure appear in \Cref{app:evaluation-timing}.

\Cref{tab:rq3-performance} summarizes the resulting timing measurements for
the default double-word wrappers and the transformation wrappers.  The
underlying times and results for the strong wrappers appear in
\cref{tab:rq3-performance-detail}.

\begin{table}[t]
  \caption{RQ3: Performance of the default double-word wrappers and the
  transformation wrappers.  Under CompCert, Wrapper/Kernel is the ratio to the
  unchecked kernel and Added is the paired difference in nanoseconds per call.
  Under Clang, Wrapper/QD is the ratio to QD.  Entries are medians
  [interquartile ranges] over paired trials.}
  \label{tab:rq3-performance}
  \centering
  \setlength{\tabcolsep}{3pt}
  \begin{tabular}{@{}lrrr@{}}
    \toprule
    & \multicolumn{2}{c}{CompCert} & Clang \\
    \cmidrule(lr){2-3}\cmidrule(l){4-4}
    Operation & Wrapper/Kernel & Added (ns) & Wrapper/QD \\
    \midrule
    \input{generated/rq3_compact_rows}\\
    \bottomrule
  \end{tabular}
\end{table}

We rotate the measurement order within paired trials to limit ordering effects
and compute differences and ratios from the paired measurements.

\paragraph{Answer to RQ3.}
\Cref{tab:rq3-performance} and \Cref{tab:rq3-performance-detail} show that the
wrappers have a noticeable relative cost because their kernels are small, but
add only a few nanoseconds per call.  Under CompCert, the default wrappers took
\RQThreeDefaultMinimumKernelRatio{}--\RQThreeDefaultMaximumKernelRatio{} times
as long as their kernels and added
\RQThreeDefaultMinimumAddedTime{}--\RQThreeDefaultMaximumAddedTime{}
nanoseconds per call.  The strong wrappers took
\RQThreeStrongMinimumKernelRatio{}--\RQThreeStrongMaximumKernelRatio{} times
as long and added
\RQThreeStrongMinimumAddedTime{}--\RQThreeStrongMaximumAddedTime{}
nanoseconds.  They were faster than the default wrappers for four operations
under CompCert but slower for all six under Clang, showing that their relative
cost depends on the operation and compiler.  The transformation wrappers took
\RQThreeTransformationMinimumKernelRatio{}--\RQThreeTransformationMaximumKernelRatio{}
times as long as their kernels.

Under Clang, the default wrappers took
\RQThreeDefaultMinimumQDRatio{}--\RQThreeDefaultMaximumQDRatio{} times as long
as QD, the strong wrappers took
\RQThreeStrongMinimumQDRatio{}--\RQThreeStrongMaximumQDRatio{} times as long,
and the transformation wrappers took
\RQThreeTransformationMinimumQDRatio{}--\RQThreeTransformationMaximumQDRatio{}
times as long.  These measurements compare total running time and therefore
also include differences between implementations and algorithms.
\FloatBarrier
